Published phishing email detectors report 97 to 99\% accuracy, almost always
measured inside a single corpus. We ask what they are worth on mail from a
corpus they have never seen, once the emails that public corpora share with
each other are removed. Comparing nine public sources, we find that
87\% of SpamAssassin, 91\% of Ling-Spam and
36\% of Enron have a near duplicate inside the widely used Kaggle
phishing set, while an exact-match check finds 2 of the
SpamAssassin copies. Removing the duplicates lowers a transfer score from
0.966 to 0.672 F1, and removing the same number of
training emails at random instead changes nothing, which separates leakage from
lost data. An LLM annotator finds that only 15 to 25\% of
the positive emails in these corpora are phishing rather than bulk spam. On
decontaminated leave-one-corpus-out data we compare classical models, a
fine-tuned DistilBERT, six small open LLMs, three published phishing detectors
and two commercial models on identical emails, with false alarm rates and the
measured dollar cost. Zero-shot Qwen-2.5-7B reaches 0.928 mean F1
with 3\% false alarms at three US cents per 1{,}000 emails,
Gemini-3.1-Flash-Lite reaches 0.960, and a two-stage detector that
sends only flagged mail to an LLM reaches 0.911 at
1.1\% false alarms for 0.0161 dollars per
1{,}000. We also show that the AI-written phishing corpus used here carries a
give-away token which a naive comparison blames for 16 points of recall,
but which a paired test clears (0.9 points), and that
few-shot examples taken from old corpora, rather than few-shot prompting
itself, are what lower performance on AI-written mail. The whole study ran on
one laptop for 1.58 US dollars.